\documentclass[twocolumn]{iopjournal}

\usepackage{amsmath,amssymb,amsfonts}
\usepackage{amsthm}

\fancyhfoffset[L]{0mm}
\setlength{\marginparwidth}{0pt}
\setlength{\marginparsep}{0pt}

% ==========================================
% ZENODO PREPRINT HEADERS (ACTIVE)
% ==========================================
\renewcommand{\articletype}[1]{%
  {\vspace*{-8mm}\noindent \Large \sf Preprint}\par
  \vspace*{6mm}\noindent{\scriptsize \sf{\bfseries \MakeUppercase{#1}}}}

\fancyhead[L]{}
\fancyhead[R]{}

% Theorem environments
\newtheorem{lemma}{Lemma}
\newtheorem{theorem}[lemma]{Theorem}
\newtheorem{corollary}[lemma]{Corollary}
\theoremstyle{definition}
\newtheorem{definition}[lemma]{Definition}
\newtheorem{remark}[lemma]{Remark}

\begin{document}

\articletype{Supplementary Material}

\title{The Constitutive Law from Exchange Symmetry}

\author{Stephen Nelson}

\affil{Field Dynamics Research}

\ead{stephen@fielddynamics.org}

\begin{abstract}
The tetrad equation produces two parity-related bounded regions,
TetA and TetB, whose spatial geometry is the stellated
octahedron (Supplementary Proof~A).  Any constitutive response
$\mu(x)$ measuring the TetA fraction at acceleration $x$ must
satisfy a duality: exchanging TetA and TetB sends $x \to 1/x$
and $\mu \to 1 - \mu$.  Defining the odds ratio
$\Omega = \mu/(1 - \mu)$, this duality becomes
$\Omega(1/x) = 1/\Omega(x)$.  The inversion $x \mapsto 1/x$
exchanges $0$ and $\infty$, so the natural compactified domain
is the projective line $\mathbb{P}^1$.  We require $\Omega$ to
be an automorphism of this symmetry space; by the standard
classification $\mathrm{Aut}(\mathbb{P}^1) =
\mathrm{PGL}_2$~\cite{Hartshorne1977}, every such map is
fractional linear.  Within this class, the exchange symmetry
and the endpoint condition $\mu(0) = 0$ together force
$\Omega(x) = x$, from which $\mu(x) = x/(x+1)$ follows by
algebraic normalisation.  The midpoint value $\mu(1) = 1/2$ is
not assumed; it is a derived consequence.
\end{abstract}

% ==================================================================

\section*{The exchange symmetry}

The tetrad equation
$g_{\mu\nu} = \eta_{ab}\, e^a{}_\mu\, e^b{}_\nu$ is quadratic
in the frame field~$e$.  In the spatial sector, it therefore
admits two roots related by spatial parity
$P: e^i \to -e^i$.  The accompanying Supplementary Proof~A
(\textit{The Spatial Geometry of the Lorentz Tetrad})
establishes that each root realises a regular tetrahedron ---
the identity component TetA and the parity component TetB ---
and that their compound is the stellated
octahedron~\cite{Kepler1619}.

\begin{definition}[Constitutive response]
\label{def:response}
Let $x = g/a_0$ be the dimensionless acceleration ratio.
Define $\mu(x) \in [0,1]$ as the fraction of the total field
capacity carried by TetA at acceleration $x$, so that
$1 - \mu(x)$ is the fraction carried by TetB.
\end{definition}

\noindent
Spatial parity exchanges TetA and TetB.  Under this exchange,
the acceleration ratio inverts ($x \to 1/x$) and the two
fractions swap:

\begin{definition}[Exchange symmetry]
\label{def:exchange}
The constitutive response satisfies the duality condition
\begin{equation}
  \mu(1/x) = 1 - \mu(x).
  \label{eq:duality}
\end{equation}
Two endpoint conditions follow from the definitions of the
TetA and TetB fractions:
$\mu(0) = 0$ (at zero acceleration, TetA is empty) and
$\mu(x) \to 1$ as $x \to \infty$ (at large acceleration,
TetA saturates).
\end{definition}

% ==================================================================

\section*{The odds ratio}

\begin{definition}[Odds ratio]
\label{def:odds}
The odds ratio between TetA and TetB is
\begin{equation}
  \Omega(x) = \frac{\mu(x)}{1 - \mu(x)}.
  \label{eq:odds}
\end{equation}
This measures how many times more capacity TetA carries than
TetB at acceleration $x$.
\end{definition}

\begin{lemma}[Odds-ratio duality]
\label{lem:odds_duality}
Under the exchange symmetry~\eqref{eq:duality}, the odds ratio
satisfies $\Omega(1/x) = 1/\Omega(x)$.
\end{lemma}

\begin{proof}
\begin{align}
  \Omega(1/x)
  &= \frac{\mu(1/x)}{1 - \mu(1/x)}
   = \frac{1 - \mu(x)}{\mu(x)} \notag \\
  &= \frac{1}{\Omega(x)}.
\end{align}
\hfill$\square$
\end{proof}

\noindent
The endpoint conditions translate to:
\begin{equation}
  \Omega(0) = 0, \qquad \Omega(x) \to \infty
  \text{ as } x \to \infty.
  \label{eq:odds_endpoints}
\end{equation}

\begin{lemma}[Symmetry domain]
\label{lem:domain}
The exchange $x \mapsto 1/x$ sends $0 \mapsto \infty$ and
$\infty \mapsto 0$.  The natural compactified domain on which
this involution acts is the projective line $\mathbb{P}^1$.
\end{lemma}

\begin{definition}[Structural axiom]
\label{def:axiom}
We require the odds map $\Omega$ to be an automorphism of the
algebraic symmetry domain $\mathbb{P}^1$.  That is,
$\Omega \in \mathrm{Aut}(\mathbb{P}^1)$.  This is the minimal
algebraically closed class containing the exchange involution
$x \mapsto 1/x$ and stable under composition and inversion.
\end{definition}

\begin{lemma}[Classification of $\mathrm{Aut}(\mathbb{P}^1)$]
\label{lem:classification}
Over $\mathbb{R}$,
$\mathrm{Aut}(\mathbb{P}^1) = \mathrm{PGL}_2(\mathbb{R})$.
Every such automorphism is fractional
linear~\cite{Hartshorne1977}:
\begin{equation}
  \Omega(x) = \frac{ax + b}{cx + d}, \qquad ad - bc \neq 0.
  \label{eq:mobius_general}
\end{equation}
\end{lemma}

% ==================================================================

\section*{Main result}

\begin{theorem}[Forced odds ratio]
\label{thm:forced}
Let $\Omega: \mathbb{P}^1 \to \mathbb{P}^1$ be a rational
bijection satisfying $\Omega(1/x) = 1/\Omega(x)$ and
$\Omega(0) = 0$.  Then
\begin{equation}
  \boxed{\Omega(x) = x.}
  \label{eq:odds_result}
\end{equation}
\end{theorem}

\begin{proof}
By Lemma~\ref{lem:classification},
$\Omega(x) = (ax+b)/(cx+d)$.
The endpoint condition $\Omega(0) = b/d = 0$ requires $b = 0$.
The map reduces to $\Omega(x) = ax/(cx + d)$.

Substituting into the exchange condition
$\Omega(1/x) = 1/\Omega(x)$:
\begin{equation}
  \frac{a}{c + dx} = \frac{cx + d}{ax}.
  \label{eq:exchange_sub}
\end{equation}
Cross-multiplying: $a^2 x = (c + dx)(cx + d)
= cd\,x^2 + (c^2 + d^2)\,x + cd$.  Comparing coefficients:
\begin{equation}
  cd = 0, \qquad a^2 = c^2 + d^2.
  \label{eq:coefficients}
\end{equation}
From $cd = 0$, either $c = 0$ or $d = 0$.

If $d = 0$: $\Omega(x) = ax/(cx) = a/c$, a constant.  This is
not bijective.  Rejected.

If $c = 0$: $\Omega(x) = (a/d)\,x$, and
equation~\eqref{eq:coefficients} gives $(a/d)^2 = 1$.  Since
$\Omega$ is an odds ratio, $\Omega(x) \geq 0$ for $x \geq 0$,
so $a/d = +1$.
\hfill$\square$
\end{proof}

\noindent
The odds ratio between TetA and TetB is the acceleration ratio
itself.  This is not assumed; it is the unique element of
$\mathrm{Aut}(\mathbb{P}^1)$ satisfying the exchange symmetry
and the endpoint condition.  The result is a uniqueness theorem
within the automorphism group of the projective symmetry space,
not a claim about arbitrary constitutive functions.

% ==================================================================

\section*{Corollaries}

\begin{corollary}[Constitutive law]
\label{cor:mu}
From $\Omega(x) = x$, inverting the
definition~\eqref{eq:odds}:
\begin{equation}
  \boxed{\mu(x) = \frac{x}{x + 1}.}
  \label{eq:mu}
\end{equation}
\end{corollary}

\begin{proof}
$\mu = \Omega/(1 + \Omega) = x/(1 + x)$.
\hfill$\square$
\end{proof}

\begin{corollary}[Midpoint value]
\label{cor:midpoint}
$\mu(1) = 1/2$.  At the parity midpoint, TetA and TetB
contribute equally.
\end{corollary}

\begin{proof}
$\mu(1) = 1/(1 + 1) = 1/2$.  This is a consequence of the
exchange symmetry, not an independent condition.
\hfill$\square$
\end{proof}

% ==================================================================

\section*{Remarks}

\begin{remark}[Complement symmetry is the starting point]
\label{rem:complement}
The identity $\mu(x) + \mu(1/x) = 1$ was not assumed as a
separate condition.  It is the duality
condition~\eqref{eq:duality}, which was the starting point of
the derivation.
\end{remark}

\begin{remark}[Screening identity is algebraic]
\label{rem:screening}
The screening identity $x(1 - \mu) = \mu$ is not imposed.  It
is the odds-ratio identification $\mu/(1-\mu) = x$ rearranged:
multiply both sides by $(1-\mu)$ to obtain $\mu = x(1-\mu)$.
This is an algebraic identity of equation~\eqref{eq:mu}, not an
independent physical postulate.
\end{remark}

\begin{remark}[What this derivation does not use]
\label{rem:independence}
This derivation uses no coupling polynomial, no rank counting,
no field equation, and no equilibrium postulates.  It requires
only the exchange symmetry between the two parity-related roots
of the tetrad equation and the classification of
$\mathrm{Aut}(\mathbb{P}^1)$.  The function
$\mu(x) = x/(x+1)$ was identified empirically by
Milgrom~\cite{Milgrom1983} and selected by fitting by Famaey
and Binney~\cite{FamaeyBinney2005}.  The present derivation
obtains it as a uniqueness theorem: given equivariance within
the automorphism group of the projective symmetry space, no
other constitutive law is possible.
\end{remark}

% ==================================================================

\begin{thebibliography}{9}

\bibitem{Kepler1619}
J.~Kepler,
\textit{Harmonices Mundi}
(Linz, 1619).

\bibitem{Hartshorne1977}
R.~Hartshorne,
\textit{Algebraic Geometry}
(Springer, New York, 1977).

\bibitem{Milgrom1983}
M.~Milgrom,
\textit{A modification of the Newtonian dynamics as a possible
alternative to the hidden mass hypothesis},
Astrophys.\ J.\ \textbf{270} (1983) 365.

\bibitem{FamaeyBinney2005}
B.~Famaey and J.~Binney,
\textit{Modified Newtonian dynamics in the Milky Way},
Mon.\ Not.\ R.\ Astron.\ Soc.\ \textbf{363} (2005) 603.

\end{thebibliography}

\end{document}
